%% main.tex — DRAFT scaffold for the volumetric-cloud 3DGS paper.
%%
%% Compile from this directory (acmart.cls sits alongside):
%%   pdflatex main && bibtex main && pdflatex main && pdflatex main
%% For submission/review switch to:
%%   \documentclass[manuscript,screen,review]{acmart}
%%
%% LEGEND for the author:
%%   [VERIFIED]  — claim/number grounded in the current code or training logs.
%%   TODO        — placeholder to fill (figures, citations, ablation numbers).
%% Search the file for "TODO" before submitting.

\documentclass[sigconf]{acmart}

\AtBeginDocument{%
  \providecommand\BibTeX{{%
    Bib\TeX}}}

\setcopyright{none}            % TODO: set real rights when known
\acmConference[Preprint]{Preprint — work in progress}{2026}{}
\settopmatter{printacmref=false}
\renewcommand\footnotetextcopyrightpermission[1]{}
\pagestyle{plain}

%% --- handy macros --------------------------------------------------------
\usepackage{amsmath}
\newcommand{\bpeak}{\beta_{\mathrm{peak}}}
\newcommand{\Tlight}{T_{\mathrm{light}}}
\newcommand{\todo}[1]{\textcolor{red}{\textbf{[TODO: #1]}}}
\usepackage{xcolor}

\begin{document}

\title{Physically-Parameterized 3D Gaussian Splatting for\\
       Real-Time Relightable Volumetric Clouds}

%% TODO: real author block / affiliations.
\author{Anonymous}
\affiliation{%
  \institution{TODO Institution}
  \country{}
}

\begin{abstract}
Real-time volumetric clouds in game engines are traditionally rendered with
ray marching, which is expensive and couples shading cost to volume density.
We instead represent a cloud as a set of 3D Gaussians whose appearance is
\emph{physically} parameterized as a participating medium---each Gaussian
carries a peak extinction coefficient, a scattering albedo, and a
Henyey--Greenstein anisotropy parameter rather than spherical-harmonic colour
and an opacity scalar. Radiative transfer is evaluated in closed form: we
derive an analytic optical-depth line integral for each Gaussian along both the
view ray and the sun ray, composite it as $\alpha = 1 - e^{-\tau}$, and pass
the per-Gaussian optical depth into a modified differentiable rasterizer,
avoiding ray marching entirely. Multiple scattering is approximated by a
six-octave Henyey--Greenstein expansion, and per-Gaussian sun transmittance is
obtained from a light-space optical-depth field. Because the sun direction is
supplied per frame and trained over a range of directions, the learned cloud
can be \emph{relit interactively} at inference. We further introduce a
contribution-aware densification/pruning strategy tailored to the physical
parameterization. On a dynamically-lit synthetic cloud dataset our method
reaches \todo{PSNR/SSIM/LPIPS final numbers} while rendering in real time.
\end{abstract}

%% TODO: pick real CCS concepts at https://dl.acm.org/ccs
\begin{CCSXML}
<ccs2012>
 <concept>
  <concept_id>10010147.10010371.10010396</concept_id>
  <concept_desc>Computing methodologies~Rendering</concept_desc>
  <concept_significance>500</concept_significance>
 </concept>
</ccs2012>
\end{CCSXML}
\ccsdesc[500]{Computing methodologies~Rendering}

\keywords{3D Gaussian Splatting, volumetric clouds, participating media,
real-time rendering, relighting, radiative transfer}

\maketitle

%% =======================================================================
\section{Introduction}
\label{sec:intro}

Volumetric clouds are a staple of real-time graphics, but the dominant
technique---ray marching a density field---spends compute proportional to the
number of samples along each ray, independent of how much the cloud actually
contributes to the image. 3D Gaussian Splatting (3DGS)~\cite{kerbl3Dgaussians}
has shown that an explicit set of anisotropic Gaussians, rasterized with
$\alpha$-compositing, can reconstruct radiance fields at real-time rates with
quality competitive with NeRF~\cite{mildenhall2020nerf}. Its primitive---a
smooth, anisotropic blob---is also a natural match for the soft, wispy
structure of clouds.

Vanilla 3DGS, however, is a \emph{view-dependent appearance} model: each
Gaussian stores spherical-harmonic colour and a free opacity scalar, baking
the lighting of the training images into the primitives. This makes relighting
impossible and discards the physics that makes a cloud a cloud. Our goal is a
representation that (i) renders in real time without ray marching, (ii) is
\emph{relightable}---the sun can be moved at inference---and (iii) exposes
physically meaningful parameters (extinction, albedo, phase anisotropy).

We replace the appearance model of 3DGS with a participating-medium model and
keep the rasterizer as the integration engine. Our contributions are:

\begin{itemize}
  \item A \textbf{physical reparameterization} of each Gaussian as a medium
        kernel: peak extinction $\bpeak$, albedo $\rho$, and
        Henyey--Greenstein anisotropy $g$, with derived optical mass
        (Sec.~\ref{sec:param}).
  \item An \textbf{analytic optical-depth rasterization}: a closed-form 1D
        Gaussian line integral yields per-Gaussian optical depth along the view
        and sun rays; we composite $\alpha=1-e^{-\tau}$ and feed $\tau$ into a
        modified differentiable rasterizer, removing ray marching
        (Sec.~\ref{sec:render}).
  \item A \textbf{six-octave Henyey--Greenstein} multiple-scattering
        approximation and a \textbf{light-space self-shadow} estimate
        $\Tlight$ that is differentiable in Gaussian position and optical depth
        (Sec.~\ref{sec:ms}, \ref{sec:shadow}).
  \item A \textbf{contribution-aware densification/pruning} scheme suited to
        the physical parameterization, and \textbf{interactive relighting} from
        per-frame sun supervision (Sec.~\ref{sec:densify}, \ref{sec:relight}).
\end{itemize}

%% =======================================================================
\section{Related Work}
\label{sec:related}

\paragraph{Radiance-field reconstruction.}
NeRF~\cite{mildenhall2020nerf} and its many successors model a scene as a
volumetric radiance field queried by ray marching. 3DGS~\cite{kerbl3Dgaussians}
replaces the implicit field with explicit anisotropic Gaussians and a tile-based
rasterizer, achieving real-time rendering. \todo{cite 2--3 follow-ups relevant
to your angle, e.g. relightable / inverse-rendering GS.}

\paragraph{Relightable Gaussian Splatting.}
\todo{Discuss GS-IR, Relightable-3D-Gaussians, and any cloud/medium GS work;
contrast their per-Gaussian visibility (depth-cubemap / BVH ray tracing /
shadow-ray opacity) with our light-space analytic optical-depth readback.}

\paragraph{Volumetric clouds and participating media.}
The Henyey--Greenstein phase function~\cite{henyey1941diffuse} is the standard
single-lobe scattering model for clouds. Real-time cloud systems approximate
multiple scattering with cheap multi-octave
expansions~\cite{wrenninge2013oz, hillenius2016frostbite}. \todo{verify and
finalize these references; add the production cloud refs you cited in the
design doc.}

%% =======================================================================
\section{Physical Gaussian Parameterization}
\label{sec:param}

%% [VERIFIED against scene/gaussian_model.py]
Each Gaussian $i$ retains its 3DGS geometry---centre $\mu_i$, scale $s_i \in
\mathbb{R}^3$, rotation $q_i$---but its appearance is re-parameterized as a
participating-medium kernel:
\begin{align}
  \bpeak^{(i)} &= \mathrm{softplus}(\theta^\beta_i) \le 5,
    &\text{(peak extinction, $1/\text{length}$)}\\
  \rho_i       &= \sigma(\theta^\rho_i) \in (0,1)^3,
    &\text{(RGB single-scatter albedo)}\\
  g_i          &= 0.8\,\tanh(\theta^g_i) \in (-0.8, 0.8).
    &\text{(HG anisotropy)}
\end{align}
The spherical-harmonic colour and free opacity scalar of vanilla 3DGS are
removed entirely. $\bpeak$ is the \emph{intensive} peak extinction, not total
mass; the optical mass of a normalized 3D Gaussian is recovered analytically as
\begin{equation}
  m_i = \bpeak^{(i)} \,(2\pi)^{3/2} \prod_d s_{i,d}.
  \label{eq:mass}
\end{equation}
Because $\bpeak$ is intensive (a local density), densification by splitting or
cloning \emph{inherits} it unchanged rather than dividing it---only the scale
changes---which keeps the medium's density physically consistent across
multi-scale refinement.

%% =======================================================================
\section{Analytic Optical-Depth Rasterization}
\label{sec:render}

%% [VERIFIED against gaussian_renderer/__init__.py + the modified rasterizer]
Rather than ray-march a density field, we integrate each Gaussian's extinction
in closed form. For a normalized 3D Gaussian with scales $s$ and a unit
direction $d$ expressed in the Gaussian's local frame, the full-line integral
through the centre is
\begin{equation}
  L(s,d) = \frac{1}{(2\pi)\,\prod_d s_d \,\sqrt{\sum_d (d_d/s_d)^2}},
  \label{eq:lineint}
\end{equation}
the exact 1D Gaussian integral. Multiplying by the mass~\eqref{eq:mass} gives
the per-Gaussian optical depth along the ray, e.g.\ along the view ray
$\tau^{\mathrm{view}}_i = m_i\,L(s_i, v_i)$, and the per-pixel opacity is
\begin{equation}
  \alpha_i = 1 - \exp\!\big(-\tau^{\mathrm{view}}_i\big).
\end{equation}
We pass $\tau^{\mathrm{view}}$ to the rasterizer as a dedicated
\texttt{tau\_precomp} channel; inside the CUDA kernel an
\texttt{use\_analytic\_tau} branch composites $\alpha = 1 - \exp(-\tau\cdot
G_{2D})$, where $G_{2D}$ is the projected 2D Gaussian footprint, so the
splat realizes the intended $\tau(x') \approx \tau_{\text{centre}}\,G_{2D}(x')$
approximation. The forward and backward passes are differentiable with respect
to $\mu$, $s$, $q$, and $\tau$ (hence $\bpeak$). Because the model sums optical
depth rather than $\alpha$-over compositing densities, the per-Gaussian
contribution is order-independent in $\tau$, matching the physics of a
participating medium. \todo{state the exact rasterizer fork and that the stock
3DGS path is retained.}

%% =======================================================================
\section{Multiple Scattering}
\label{sec:ms}

%% [VERIFIED: 6-octave loop in render()]
Single scattering alone makes clouds look thin and dark. We approximate
multiple scattering with a six-octave Henyey--Greenstein
expansion~\cite{wrenninge2013oz}: octave $n$ attenuates energy by $a^n$, raises
transmittance to the power $b^n$, and isotropizes the phase by $g_{\text{eff}} =
g\,c^n$, with $a=b=c=0.5$. The outgoing radiance of Gaussian $i$ is
\begin{equation}
  L_i = \rho_i\, L_{\mathrm{sun}}
        \sum_{n=0}^{5} a^n\, \Tlight^{\,b^n}\,
        \mathrm{HG}\!\big(g_i c^n, \cos\theta\big),
\end{equation}
where $\mathrm{HG}$ is the $1/4\pi$-normalized phase function and $L_{\mathrm{sun}}
= 4\pi$ compensates the normalization so an isotropic unit-albedo medium
reflects all incident light.

%% [VERIFIED: HG sign convention — this session's finding]
\paragraph{Phase-angle convention.}
The HG phase function is defined on the angle between the photon's
\emph{incoming} propagation direction $\omega_{\text{in}}$ and its outgoing
direction $\omega_{\text{out}}$. With a sun direction that points
\emph{toward} the sun, the light propagates along $\omega_{\text{in}} =
-\,\hat{\ell}$, so $\cos\theta = (-\hat{\ell})\cdot v$. Using $+\hat{\ell}$
instead flips the forward lobe and forces the optimizer to learn a spurious
\emph{negative} $g$ for a physically forward-scattering cloud; with the correct
convention the learned anisotropy settles at a physical, mildly
forward-scattering value (Sec.~\ref{sec:results}).

%% =======================================================================
\section{Light-Space Self-Shadow}
\label{sec:shadow}

%% [VERIFIED: compute_T_light()]
Each Gaussian must be shadowed by the medium between it and the sun. We build an
orthonormal light-space frame whose third axis is the sun direction, deposit
each Gaussian's sun-ray optical depth $\tau^{\mathrm{sun}}_i = m_i\,L(s_i,
\ell_i)$ onto a $128^3$ voxel grid, and take an exclusive prefix sum along the
sun axis to obtain, per voxel, the optical depth accumulated between it and the
sun. Trilinear sampling at each Gaussian centre yields $\tau^{\mathrm{sun,acc}}_i$
and the transmittance $\Tlight^{(i)} = \exp(-\tau^{\mathrm{sun,acc}}_i)$.
The readback is differentiable in the Gaussian position (through the sample
coordinate) and in the deposited optical depth (through $\bpeak$, $s$).
\todo{add the deep-opacity-map / sun-view rasterization discussion as future
work; note the voxel approximation's resolution trade-off.}

%% =======================================================================
\section{Densification and Pruning}
\label{sec:densify}

%% [VERIFIED: physical_densify_and_prune / tick_post_densify_maintenance]
The opacity-threshold pruning of vanilla 3DGS does not apply once opacity is
derived from physics. We prune on three signals: (1) \emph{contribution}---the
running mean of $\sum(\alpha\,T)$ a Gaussian accrues over the frames in which it
is visible, removing ``ghost'' Gaussians with negligible image contribution;
(2) \emph{dead points}---Gaussians visible in zero frames; and (3)
\emph{anisotropy}---Gaussians whose axis ratio exceeds a threshold, which we
find is the effective control for depth-sort popping (Sec.~\ref{sec:results}).
A per-Gaussian grace counter grants newly split/cloned/resurrected points a
fixed window of prune immunity, and densification is gated to Gaussians that are
alive or within grace, so dead points are not cloned. A gentle log-ratio
anisotropy regularizer is applied during densification and disabled afterward.

%% =======================================================================
\section{Interactive Relighting}
\label{sec:relight}

%% [VERIFIED: per-frame sun_dir through Camera -> render()]
The sun direction is supplied per training frame and threaded from the dataset
through the camera to the shader and the self-shadow estimate. Training over a
range of sun directions yields a single cloud model whose lighting is a runtime
input: at inference the sun can be dragged and $\Tlight$ recomputed, relighting
the cloud in real time. \todo{report relighting frame rate and show a sun-sweep
figure.}

%% =======================================================================
\section{Experiments}
\label{sec:results}

\paragraph{Setup.}
%% [VERIFIED: dataset reader, train.py loss]
We train on a synthetic cloud dataset exported from a game engine (UE5),
converted to the NeRF-synthetic transform convention, with per-frame sun
directions. The loss is $(1-\lambda)\,\mathcal{L}_1 + \lambda\,(1-\text{SSIM})$
plus a volume regularizer $\lambda_{\text{vol}}\,\overline{\prod_d s_d}$ and the
anisotropy term. We optimize for 30k iterations. \todo{exact split sizes,
resolution, hardware, frame rate, optimizer LRs.}

\paragraph{Reconstruction quality.}
%% [VERIFIED numbers: latest dynamic-dataset run @30k]
On the dynamic dataset our final configuration reaches train
PSNR~$\approx 36.0$, SSIM~$\approx 0.986$, LPIPS~$\approx 0.067$
(Table~\ref{tab:main}). \todo{replace train metrics with held-out test metrics
and add baselines.}

\begin{table}[t]
  \caption{Reconstruction on the dynamic cloud dataset. \todo{add test split,
  baselines, and a real frame-rate column.}}
  \label{tab:main}
  \begin{tabular}{lccc}
    \toprule
    Configuration & PSNR\,$\uparrow$ & SSIM\,$\uparrow$ & LPIPS\,$\downarrow$\\
    \midrule
    Full model (ours) & 36.01 & 0.986 & 0.068 \\
    \todo{baseline}    & --    & --    & --    \\
    \bottomrule
  \end{tabular}
\end{table}

\paragraph{Phase-anisotropy sign (ablation).}
%% [VERIFIED: this session]
With the incorrect phase-angle sign the optimizer learns $g \approx -0.14$, a
non-physical apparent back-scattering, at PSNR~$35.9$. Correcting the sign
yields $g \approx +0.24$ (mild forward scattering, physical) at essentially
unchanged PSNR ($35.8$--$36.0$), confirming the sign is a relabeling that
recovers physical interpretability without trading quality. \todo{add the
global-vs-per-point $g$ probe and a phase-function plot; discuss why the learned
$g$ is lower than a textbook single-scatter cloud value (multi-octave
isotropization dilutes the nominal $g$).}

\paragraph{Popping control (ablation).}
%% [VERIFIED: k_sigma revert this session]
A per-tile max-response depth sort was tried to fix long-axis popping of
elongated Gaussians but produced blocky tile-boundary artifacts; the anisotropy
prune/penalty controls popping on its own, so we revert to the stock
centre-depth sort. This both removes the artifacts and leaves PSNR unchanged.
\todo{side-by-side popping/artifact figure.}

%% =======================================================================
\section{Limitations and Future Work}
\label{sec:limits}

The self-shadow estimate is a fixed-resolution voxel grid; a light-space
optical-depth rasterization (a differentiable deep-opacity map reusing our own
analytic-$\tau$ rasterizer) would remove the grid quantization and is a natural
next step. \todo{note any remaining capacity ceiling on the dynamic dataset and
the multi-octave $g$-dilution caveat.}

%% =======================================================================
\section{Conclusion}
\label{sec:conclusion}

We presented a physically-parameterized 3D Gaussian Splatting model for
volumetric clouds that renders in real time via analytic optical-depth
rasterization, supports interactive relighting from per-frame sun supervision,
and exposes physically meaningful medium parameters. \todo{one-sentence
quantitative takeaway once final numbers are in.}

%% =======================================================================
\bibliographystyle{ACM-Reference-Format}
\bibliography{refs}

\end{document}
